% ===================================================================
%  BẢNG Số 14 — Đường phố. --- Người buôn bán.
%
%  Traduction, faite en 2026, en vietnamien standard d'aujourd'hui,
%  sur l'ido de texto/io. Regles de la colonne : voir 10-bang-01.tex.
%
%  QUARANTE METIERS EN UN TABLEAU, ET LE VIETNAMIEN LES RANGE EN
%  TROIS COLONNES SANS QU'ON AIT A LE DECIDER. Le classificateur fait
%  seul le travail, comme au tableau 5 :
%
%    thợ    celui qui fabrique      thợ ảnh le photographe, thợ cắt
%                                   tóc le coiffeur, thợ đóng sách le
%                                   relieur, thợ nguội le mecanicien
%    người  celui qui fait l'action  người quét đường le balayeur,
%                                   người bán hàng rong le colporteur
%    hàng   celui qui vend           hàng mũ le chapelier, hàng bánh
%                                   le patissier, hàng hoa la
%                                   fleuriste, hàng thuốc lá le
%                                   bureau de tabac
%
%  « Hàng » nomme a la fois LA BOUTIQUE ET CELUI QUI LA TIENT, et le
%  livret en profite : la ou le francais devait ecrire « le magasin du
%  chapelier », le vietnamien ecrit « hàng mũ » et laisse le lecteur
%  voir les deux.
%
%  DEUX PIEGES DU QUEBEC N'EN SONT PAS ICI, ET UN TROISIEME LE
%  DEVIENT. La balayeuse (12) ne se confond avec aucun aspirateur —
%  « xe quét đường », la voiture qui balaie la rue. La borne-fontaine
%  (78) ne se confond avec aucune bouche d'incendie — « vòi nước
%  công cộng », le robinet public. Mais l'avertisseur d'incendie (76)
%  et le robinet public sont voisins de trois lignes dans l'alinea, et
%  les deux commencent par « vòi » : on ecrit donc le premier « chuông
%  báo cháy », la cloche qui annonce le feu, pour qu'aucun des deux ne
%  prenne le nom de l'autre.
%
%  LA MACHINE A ECRIRE EST « MAY CHU », la machine a caracteres, et
%  c'est un mot de 1926 : elle etait alors la nouveaute meme du
%  bureau, et le vietnamien l'a nommee sans l'emprunter.
% ===================================================================

%%K t14-apar-1 apar
\VUsaut{4.0mm}
\VUtitre{15.0pt}{60}{BẢNG Số 14}
\VUsaut{3.0mm}

%%K t14-tit-1 sub
\VUpk{11.4pt}{40}{Đường phố. --- Người buôn bán.}
\VUsaut{3.0mm}

%%K t14-01-1 p
1. --- Bạn em, Gioan, ở nhà quê, đến chơi nhà em ba ngày. \VUgras{Cho đến lúc ấy} bạn ấy
chưa có dịp thấy một thành phố lớn; vì thế chúng em dùng cả ngày để đi dạo, và em giảng
cho bạn những gì bạn chưa biết.

%%K t14-02-1 p
2. --- Trước hết chúng em đi thăm \VUgras{đài thiên văn} \textsuperscript{(1)}, nơi làm
bạn ấy rất thích thú. Khi về theo \VUgras{con phố} \textsuperscript{(3)} có \VUgras{nhà
tắm hơi kiểu Thổ} \textsuperscript{(2)}, chúng em gặp một \VUgras{đám ma}
\textsuperscript{(4)}. Đi trước \VUgras{đoàn đưa tang} là \VUgras{hàng giáo sĩ}, dẫn
trước \VUgras{xe tang} chở chiếc \VUgras{quan tài}. Rất nhiều bạn bè đi cùng gia đình
\VUgras{để tang}.

%%K t14-02-2 p
Sau khi đoàn đưa tang qua rồi, chúng em băng qua phố trước \VUgras{quán bia}
\textsuperscript{(5)} lớn, nơi nhiều \VUgras{khách} đang uống \VUgras{bia} trên sân
thượng, dưới bóng tấm \VUgras{mái hiên kính} \textsuperscript{(6)}.

%%K t14-03-1 p
3. --- Gioan rất ngạc nhiên về tấm \VUgras{huy hiệu} treo giữa cửa hàng của \VUgras{người
bán rượu mạnh} \textsuperscript{(7)} và cửa hàng của \VUgras{người bán nhạc cụ}
\textsuperscript{(10)}. Bạn hỏi em nó nghĩa là gì; em trả lời rằng nó chỉ \VUgras{tòa
lãnh sự} \textsuperscript{(8)} Anh, và em chỉ cho bạn thấy \VUgras{ông lãnh sự}
\textsuperscript{(9)} đang đứng trên ban công.

%%K t14-tit-2 sub
\VUpk{10.2pt}{40}{Các tủ kính.}
\VUsaut{3.0mm}

%%K t14-04-1 p
4. --- Dĩ nhiên chúng em dừng lại trước tủ bày \VUgras{nhạc cụ}, những chiếc \VUgras{đàn
hạc phong cầm}, \VUgras{đàn ống} và \VUgras{đàn dương cầm}; chúng em xem những bìa có
tranh của các \VUgras{bản nhạc} và các \VUgras{bản soạn} cho dương cầm, và thán phục cây
\VUgras{đàn lia} bằng gỗ thếp vàng dùng làm biển hiệu.

%%K t14-05-1 p
5. --- Những đám bụi do \VUgras{người quét đường} \textsuperscript{(11)} và \VUgras{xe
quét đường} \textsuperscript{(12)} bốc lên buộc chúng em phải đi nhanh qua trước những
bộ \VUgras{bàn ghế} đẹp của \VUgras{người bán đồ gỗ} \textsuperscript{(13)} và trước tủ
bày \VUgras{bếp lò} và \VUgras{đèn} của \VUgras{hàng kim khí} \textsuperscript{(14)}.

%%K t14-06-1 p
6. --- Vả lại, ở chỗ ấy chúng em buộc phải rời vỉa hè, vì nó bị chắn bởi những kiện hàng
người ta dỡ từ một chiếc \VUgras{xe dọn nhà} \textsuperscript{(16)} để đưa vào một
\VUgras{cửa hàng} \textsuperscript{(15)} vừa được thuê.

%%K t14-06-2 p
Việc ấy làm chúng em không thấy được những cỗ xe đẹp của \VUgras{người đóng xe}
\textsuperscript{(17)}, nhưng không ngăn chúng em dừng lại xem \VUgras{người đóng kiện}
\textsuperscript{(19)} đang đóng một cái thùng cho người hàng xóm của mình,
\VUgras{người bán xe đạp và ô-tô} \textsuperscript{(18)}. Sau đó, vì đã hơi muộn, chúng
em về nhà.

%%K t14-tit-3 sub
\VUpk{10.2pt}{40}{Dạo phố.}
\VUsaut{3.0mm}

%%K t14-07-1 p
7. --- Hôm sau, mẹ em đề nghị cho Gioan đi chụp ảnh, rồi giao em đưa bạn đến nhà
\VUgras{thợ ảnh} \textsuperscript{(20)}. Sau bữa trưa, chúng em đến \VUgras{xưởng} của
người nghệ sĩ ấy, và phải chờ khá lâu. Để giết thì giờ, chúng em ngắm những \VUgras{bức
chân dung} \textsuperscript{(22)} treo trên tường.

%%K t14-07-2 p
Đến lượt chúng em, công việc không lâu, và vừa khi bạn em ngồi làm mẫu xong trước chiếc
\VUgras{máy ảnh} \textsuperscript{(21)}, chúng em xuống nhà.

%%K t14-08-1 p
8. --- Ở tầng một, qua cánh cửa để ngỏ của \VUgras{thợ cắt tóc}
\textsuperscript{(23)}, chúng em thấy người làm của ông đang cắt tóc cho một người khách.

%%K t14-08-2 p
Ra đến phố, chúng em đi qua trước \VUgras{hàng thuốc lá} \textsuperscript{(24)}. Gioan
thích thú với chiếc \VUgras{đèn lồng} \textsuperscript{(25)} đỏ và cái \VUgras{tẩu to}
dùng làm \VUgras{biển hiệu} \textsuperscript{(26)} đến nỗi bạn ấy đâm sầm vào hai ông già
đang trò chuyện và \VUgras{hít thuốc bột} \textsuperscript{(27)}.

%%K t14-tit-4 sub
\VUpk{10.2pt}{40}{Hàng bánh ngọt.}
\VUsaut{3.0mm}

%%K t14-09-1 p
9. --- Những tờ \VUgras{giấy bạc} và những đồng \VUgras{tiền kim loại} của mọi nước bày
trong \VUgras{tủ kính} của \VUgras{người đổi tiền} \textsuperscript{(29)} giữ chân chúng
em một lát, nhưng vì đã đến giờ ăn xế, chúng em vào \VUgras{hàng bánh ngọt}
\textsuperscript{(31)} mà không dừng lâu hơn.

%%K t14-10-1 p
10. --- Trước tất cả những \VUgras{món quà miệng} bày trong cửa hàng — \VUgras{bánh
ngọt}, \VUgras{bánh mật ong}, \VUgras{bánh quy} và đủ loại \VUgras{kẹo} — chúng em khó
mà chọn được. Cuối cùng Gioan chọn \VUgras{bánh su kem}, còn em chỉ lấy một chiếc
\VUgras{bánh nhân anh đào} nhỏ, vì đồ ngọt \VUgras{làm em đau răng} và em chẳng muốn đi
thăm \VUgras{nha sĩ} \textsuperscript{(30)} quá thường.

%%K t14-tit-5 sub
\VUpk{10.2pt}{40}{Máy chữ.}
\VUsaut{3.0mm}

%%K t14-11-1 p
11. --- Để cho bạn em xem một chiếc \VUgras{máy chữ} mà bạn ấy đã nghe nói nhưng chưa
biết, chúng em vào \VUgras{văn phòng} \textsuperscript{(32)} của \VUgras{anh họ}
\textsuperscript{(34)} em. Chúng em thấy anh đang đọc thư cho \VUgras{người thư ký}
\textsuperscript{(35)}, một \VUgras{người tốc ký}. Vừa đọc xong một lá thư thì một
\VUgras{người đánh máy} \textsuperscript{(33)} chép lại ngay bằng máy. Vì không muốn làm
gián đoạn công việc của người bà con, chúng em chỉ ở lại chỗ anh vài phút.

%%K t14-12-1 p
12. --- Phía dưới văn phòng của anh họ em có một \VUgras{hiệu đồng hồ và nữ trang}
\textsuperscript{(37)} lớn, chủ hiệu còn là thợ kim hoàn. Cửa hàng lộng lẫy của ông chứa
nhiều \VUgras{đồng hồ}, \VUgras{đồng hồ quả lắc}, \VUgras{đồng hồ bỏ túi},
\VUgras{dây chuyền} và những món \VUgras{nữ trang} quý, chẳng hạn \VUgras{chuỗi ngọc
trai}, \VUgras{vòng tay} vàng, \VUgras{hoa tai} và \VUgras{nhẫn} nạm \VUgras{đá quý} và
\VUgras{kim cương}. Một trong các tủ kính dành riêng cho \VUgras{đồ vàng bạc} và chứa
một bộ sưu tập lớn \VUgras{bộ đồ ăn} bằng bạc.

%%K t14-13-1 p
13. --- Rẽ ở góc phố, em nhận thấy người ta đã đổi tấm \VUgras{biển}
\textsuperscript{(36)} gắn cạnh \VUgras{số nhà}. Em sắp nói to điều ấy thì nghe tiếng
nhạc vui của \VUgras{đội} quân nhạc. Chúng em chạy ra đón trung đoàn, không nghĩ rằng
\VUgras{nó} sắp đi qua trước cửa hàng của \VUgras{người bán mũ} \textsuperscript{(39)} và
\VUgras{người bán găng tay} \textsuperscript{(40)}, và rằng đứng trước tủ kính của
\VUgras{người bán kính} \textsuperscript{(38)}, đầy \VUgras{kính kẹp mũi},
\VUgras{kính đeo} và \VUgras{ống nhòm}, thì chúng em cũng xem được cuộc diễu hành chẳng
kém.

%%K t14-tit-6 sub
\VUpk{10.2pt}{40}{Cuộc diễu hành.}
\VUsaut{3.0mm}

%%K t14-14-1 p
14. --- Đi đầu \VUgras{trung đoàn} \textsuperscript{(41)} là \VUgras{người chỉ huy đội
trống} \textsuperscript{(42)}, theo sau là những người \VUgras{đánh trống}
\textsuperscript{(43)}, những người \VUgras{thổi kèn hiệu} \textsuperscript{(44)} và cả
\VUgras{ban quân nhạc}. \VUgras{Viên tướng} \textsuperscript{(45)}, đang ở trên quảng
trường cùng sĩ quan tùy tùng, đã dừng lại gần \VUgras{vòi phun}
\textsuperscript{(49)} của \VUgras{bể nước} \textsuperscript{(48)} để dự \VUgras{cuộc
duyệt binh}.

%%K t14-14-2 p
\VUgras{Các sĩ quan tham mưu} \textsuperscript{(46)}, \VUgras{viên đại tá} và
\VUgras{viên trung tá} chào ông bằng kiếm. \VUgras{Các thiếu tá} đi đầu \VUgras{tiểu
đoàn} \textsuperscript{(47)} của mình và mỗi \VUgras{đại úy} chỉ huy \VUgras{đại đội} của
mình, trong khi các \VUgras{trung úy}, \VUgras{thiếu úy} và \VUgras{hạ sĩ quan} đứng vào
chỗ thường lệ bên cạnh quân mình.

%%K t14-15-1 p
15. --- Rất nhiều người qua đường tụ lại ở \VUgras{ngã tư} \textsuperscript{(50)}; những
người buôn bán, \VUgras{người bán ô} \textsuperscript{(51)}, \VUgras{thợ nhuộm}
\textsuperscript{(53)} và \VUgras{người bán súng} \textsuperscript{(54)} ra xem
\VUgras{lính}. Mọi thứ xe cộ — \VUgras{xe ngựa thuê} \textsuperscript{(52)}, \VUgras{xe
nhà} \textsuperscript{(57)} và cả \VUgras{xe téc tưới đường} \textsuperscript{(56)} —
đều đỗ dọc vỉa hè.

%%K t14-tit-7 sub
\VUpk{10.2pt}{40}{Cảnh đường phố.}
\VUsaut{3.0mm}

%%K t14-15-2 p
Một \VUgras{người chào hàng} \textsuperscript{(55)} xách tay những \VUgras{hộp hàng mẫu}
của mình vội vã băng qua phố, và những người ngồi trên \VUgras{tầng trên}
\textsuperscript{(59)} của chiếc \VUgras{xe khách} \textsuperscript{(58)} quay lại để
xem cho rõ. Một \VUgras{người hát rong} \textsuperscript{(60)} nghèo với chiếc \VUgras{đàn
quay} \textsuperscript{(61)} của mình phải ngừng \VUgras{bài hát}, vì tiếng trống át mất
giọng ông.

%%K t14-16-1 p
16. --- Khi trung đoàn đến trước \VUgras{cửa hàng vải} \textsuperscript{(64)}, những
\VUgras{cô thợ may} \textsuperscript{(63)} và những \VUgras{thợ may}
\textsuperscript{(62)} bỏ \VUgras{máy khâu} và công việc để nhìn qua cửa sổ.
\VUgras{Người chủ hiệu} \textsuperscript{(65)}, vừa bày ra những bộ \VUgras{quần áo}
\textsuperscript{(66)} mới trên \VUgras{giá hàng} của mình, phải cho đem vào trong những
tấm \VUgras{dạ} \textsuperscript{(67)} và những tấm \VUgras{vải} bày trên vỉa hè, gần
\VUgras{miệng cống} \textsuperscript{(68)}, sợ đám đông làm đổ; và ngay cả một \VUgras{người
bán hàng rong} \textsuperscript{(69)} già, người mà một bà gác cổng đang mặc cả mua
\VUgras{tất}, cũng phải vào một hành lang để bán cho xong.

%%K t14-16-2 p
Chúng em theo trung đoàn đến tận doanh trại, \VUgras{và}, rất hài lòng về ngày hôm ấy,
chúng em về nhà đúng giờ ăn tối.

%%K t14-tit-8 sub
\VUpk{10.2pt}{40}{Các nghề nghiệp.}
\VUsaut{3.0mm}

%%K t14-17-1 p
17. --- Hôm sau, cha em đề nghị chúng em đi thăm nhà in của tờ báo thành phố. Chúng em
hăng hái nhận lời.

%%K t14-17-2 p
\VUgras{Người chủ nhà in} vừa là \VUgras{người bán sách} \textsuperscript{(70)}
\VUgras{(= người bán sách)}, vừa là \VUgras{nhà xuất bản}, \VUgras{người bán giấy} và
\VUgras{thợ đóng sách}. Ông đã đặt ở tầng một \VUgras{tòa soạn} \textsuperscript{(71)}
của tờ báo, cùng với \VUgras{xưởng khắc}, nơi làm việc của những \VUgras{thợ in đá}
\textsuperscript{(72)} và \VUgras{thợ khắc đồng} \textsuperscript{(73)}, và sau cùng
\VUgras{xưởng sắp chữ}, nơi có những \VUgras{thợ sắp chữ} \textsuperscript{(74)}. Chính
\VUgras{nhà in} \textsuperscript{(75)}, những \VUgras{máy in} và các máy móc khác thì ở
tầng trệt.

%%K t14-tit-9 sub
\VUpk{10.2pt}{40}{Gioan hỏi rất nhiều.}
\VUsaut{3.0mm}

%%K t14-18-1 p
18. --- Ra khỏi nhà in, Gioan rất ngạc nhiên dừng lại trước một cái hộp nhỏ có kính gắn
trên một cây cột, đối diện cửa hàng của \VUgras{người bán đồ khâu vá}
\textsuperscript{(77)}, và hỏi nó dùng để làm gì. Cha em giảng cho bạn ấy rằng đó là một
\VUgras{chuông báo cháy} \textsuperscript{(76)}. Vả lại, mọi thứ trong thành phố đều làm
bạn em kinh ngạc : từ chiếc \VUgras{vòi nước công cộng} \textsuperscript{(78)} bằng đá
đơn sơ và chiếc \VUgras{xe ba bánh chở hàng} \textsuperscript{(80)} nhẹ do một
\VUgras{chú bồi} \textsuperscript{(79)} mặc chế phục đạp, cho đến chiếc \VUgras{xe thư}
\textsuperscript{(81)}.

%%K t14-19-1 p
19. --- Trong lúc cha em rời chúng em một lát để nói chuyện với \VUgras{viên thu thuế}
\textsuperscript{(83)}, người đang dừng lại trò chuyện với một \VUgras{luật sư}
\textsuperscript{(82)} và một \VUgras{chưởng khế} \textsuperscript{(84)}, chúng em vui
thú đưa mắt theo dòng người đi lại trên phố. Đủ mọi hạng người đi qua trước mặt chúng em
: một \VUgras{người làm bánh} \textsuperscript{(85)}, mang trong giỏ một chiếc
\VUgras{bánh nhân thịt} tuyệt đẹp, đi sau một \VUgras{người bán quần áo cũ}
\textsuperscript{(86)}; một \VUgras{người thắp đèn đường} \textsuperscript{(87)} trò
chuyện với một \VUgras{người coi khí đốt} \textsuperscript{(88)}; một \VUgras{nữ tu}
\textsuperscript{(89)}, tay dắt một cô bé mồ côi, đang về tu viện của mình.

%%K t14-tit-10 sub
\VUpk{10.2pt}{40}{Trên đường về.}
\VUsaut{3.0mm}

%%K t14-20-1 p
20. --- Đúng lúc cha em trở lại với chúng em, Gioan cười phá lên khi thấy khuôn mặt lem
luốc và bộ dạng kỳ quặc của hai \VUgras{người nạo ống khói} \textsuperscript{(91)} đầy
bồ hóng.

%%K t14-21-1 p
21. --- Em muốn mua của \VUgras{người bán hoa} \textsuperscript{(92)} một chậu cây cho
mẹ, nhưng cha em bảo rằng nó sẽ nặng lắm, và tốt hơn nên mua \VUgras{cam}. Chúng em lại
gần \VUgras{bà hàng} \textsuperscript{(94)} và, khi \VUgras{cô gia sư}
\textsuperscript{(93)} đang chọn cho học trò của mình mua xong, chúng em được phục vụ.

%%K t14-22-1 p
22. --- Trên đường về nhà, chúng em gặp mấy người quen : một bà bạn già của gia đình,
tựa tay vào \VUgras{người bạn cùng đi} \textsuperscript{(95)} của bà, ông \VUgras{kỹ sư}
\textsuperscript{(96)} cầu đường, và một chị họ của em cùng \VUgras{người giữ trẻ}
\textsuperscript{(98)} của chị.

%%K t14-22-2 p
Rồi chúng em gặp \VUgras{bà thợ mũ} \textsuperscript{(97)} của mẹ em và hai \VUgras{nhà
tu} \textsuperscript{(99)} lạ đến thành phố, họ lại gần hỏi cha em đường ra ga.
\VUsaut{6.0mm}
\VUfilet{22mm}
